\section{Theta-stratifications (Yaoxiong Wen and Bin Wang)}%
\label{sec:Theta-stratifications}

\subsection{Our first moduli problems}%
\label{sub:Our first moduli problem}

Our first moduli space will be \(\P^{N-1}\), which parameterizes \(1\)-dimensional linear subspaces in \(\C^N\), or equivalently is the geometric quotient \((\C^N \setminus \{0\}) / \C^{\times}\). In this case every \(\C^{\times}\)-orbit in \(\C^N\setminus\{0\}\) is closed. This construction is our first example of geometric invariant theory. In general, a set-valued moduli problem assigns to every scheme \(T\) a collection of isomorphism classes of families over \(T\), and we can think of it as a functor
\[ F \colon (\ms{Sch}_{/S})^{\mr{op}} \to \ms{Set}. \]

\begin{exm}
    Let \(C\) be a smooth projective curve. Then consider the assignment
    \[
    \begin{aligned}
        F_{r,d}\colon T \longmapsto
        \Bigl\{
        E \ \Bigm|\ &
        E \text{ is a vector bundle on } C\times T,\\
        &\operatorname{rk}(E_{\bar t})=r,\quad \deg(E_{\bar t})=d
        \text{ for every geometric point } \bar t\to T
        \Bigr\}\big/\!\sim,
    \end{aligned}
    \]
    where \(E_{\bar t}=E|_{C\times\{\bar t\}}\) and \(\sim\) denotes isomorphism of vector bundles on \(C\times T\). This parameterizes isomorphism classes of families of vector bundles on \(C\), but it is still not the correct stack-valued moduli problem.
\end{exm}

\begin{exm}
    By the Yoneda Lemma, every scheme \(Y\) defines a moduli problem.
\end{exm}

Unfortunately, not every groupoid-valued moduli problem can be represented, as a stack, by a scheme or algebraic space. A scheme or algebraic space has trivial inertia, whereas objects in a moduli problem may have nontrivial automorphisms. For example, if we consider the moduli of line bundles on \(\P^1\), each line bundle has an entire \(\C^{\times}\) of automorphisms. Therefore, when we consider moduli problems, we need to remember both the objects and the isomorphisms between them, and the target of the functor should be the category of groupoids.

\begin{defn}[informal]
    A \textit{stack} \(\mf{X}\) over \(S\) is a pseudofunctor
    \[ (\ms{Sch}_{/S})^{\mr{op}} \to \ms{Grpd} \]
    which satisfies descent for objects and morphisms with respect to the fppf topology.
\end{defn}

\begin{exm}
    Consider the stack \(\ms{Bun}_{r,d}(C)\) which sends a scheme \(T\) to the groupoid whose objects are vector bundles \(E\) on \(T\times C\) such that every geometric fiber \(E_{\bar t}\) has rank \(r\) and degree \(d\), and whose morphisms are isomorphisms of vector bundles on \(T\times C\).
\end{exm}

Unfortunately, there are far too many stacks to consider, so we will need to consider a subclass of stacks which is amenable to the methods of algebraic geometry.

\begin{defn}
    An \textit{algebraic stack} is a stack with representable (by algebraic spaces) diagonal and with a smooth surjective morphism from a scheme \(U\) (which is called an atlas). 
\end{defn}

Note that the first condition means that for any morphism \(T \to \mf{X} \times \mf{X}\) from an algebraic space, then \(T \times_{\mf{X} \times \mf{X}} \mf{X}\) is an algebraic space. This is equivalent to all Isom sheaves being algebraic spaces and implies that any morphism from a scheme to \(\mf{X}\) is representable.

\begin{rmk}
    For sets, if we have morphisms \(f\colon A\to C\) and \(g\colon B\to C\), then the fiber product is
    \[ 
        A\times_C B=\ab\{(a,b)\mid f(a)=g(b)\}. 
    \] 
    For schemes or algebraic spaces, the same formula describes the \(T\)-points of the scheme-theoretic fiber product. For stacks, a \(T\)-point of the \(2\)-fiber product \(\mf{A}\times_{\mf{C}}\mf{B}\) is a triple \((a,b,\phi)\), where \(a\in\mf{A}(T)\), \(b\in\mf{B}(T)\), and
    \[ 
        \phi\colon f(a)\xrightarrow{\sim}g(b) 
    \]
    is an isomorphism in \(\mf{C}(T)\). Thus equality must be replaced by an isomorphism, and the isomorphism must be recorded as part of the data.
\end{rmk}


\subsection{Quotient stacks}%
\label{sub:Quotient stacks}

Let \(G\) be an algebraic group and let \(X\) be a scheme with a left action of \(G\).

\begin{defn}
    The \textit{quotient stack} \([X/G]\) sends a scheme \(T\) to the groupoid whose objects are diagrams
    \begin{equation*}
    \begin{tikzcd}
        P \ar{r}{\varphi} \ar{d} & X \\
        T,
    \end{tikzcd}
    \end{equation*}
    where \(P\to T\) is a right principal \(G\)-bundle for the fppf topology and \(\varphi\) satisfies \(\varphi(pg)=g^{-1}\varphi(p)\). A morphism from \((P,\varphi)\) to \((P',\varphi')\) is an isomorphism of principal \(G\)-bundles \(P\xrightarrow{\sim}P'\) compatible with the maps to \(X\).
\end{defn}

\begin{exm}
    Let \(\mf{X}\) be a smooth separated DM stack of finite type over \(\C\). If either the generic stabilizer of \(\mf{X}\) is trivial or the coarse moduli space of \(\mf{X}\) is quasi-projective, then \(\mf{X}\) is a global quotient~\cite{EHKV2001}. Here, a DM stack is an algebraic stack with unramified diagonal, and a global quotient means a stack of the form \([U/\on{GL}_n]\) for some algebraic space \(U\) with an action of \(\on{GL}_n\).
\end{exm}

\begin{exer}
    Show that the groupoid \([X/G](\C)\) is equivalent to the action groupoid of \(G(\C)\) acting on \(X(\C)\).
\end{exer}

After choosing one representative \(x_i\in X(\C)\) of each \(G(\C)\)-orbit, the above exercise implies that, as an abstract groupoid and not as an algebraic stack, we have
\[ 
    [X/G](\C) \simeq \bigsqcup_{[x_i] \in X(\C)/G(\C)} BG_{x_i}, 
\]
where \(G_{x_i}\) is the stabilizer of \(x_i\).


\subsection{Fine vs coarse vs good moduli spaces}%
\label{sub:Fine vs coarse vs good moduli spaces}

\begin{defn}
    A \textit{fine moduli space} for a moduli problem is a scheme or algebraic space \(M\) which represents the functor. 
\end{defn}

Unfortunately, it is very hard to find a fine moduli space, so we want to weaken this condition.

\begin{defn}
    A \textit{coarse moduli space} for an algebraic stack \(\mf{X}\) is a morphism \(\pi\colon\mf{X}\to M\) to an algebraic space such that
    \begin{enumerate}
        \item For every algebraic space \(Y\), composition with \(\pi\) induces a bijection
        \[
            \on{Hom}(M,Y)\xrightarrow{\sim}\on{Hom}(\mf{X},Y);
        \]
        \item For every algebraically closed field \(K\), the induced map
        \[
            \pi_0(\mf{X}(K))=\ab\{\text{isomorphism classes in }\mf{X}(K)\}
            \longrightarrow M(K)
        \]
        is a bijection.
    \end{enumerate}
\end{defn}


Even these do not always exist. A good moduli space is a different weakening: it replaces the requirement of a bijection on geometric points by a cohomological affineness condition~\cite{Alper2013}.

\begin{defn}
    A quasi-compact morphism \(q\colon\mf{X}\to X\) from an algebraic stack to an algebraic space is a \textit{good moduli space} if
    \begin{enumerate}
        \item The pushforward functor \(q_*\colon\ms{QCoh}(\mf{X})\to\ms{QCoh}(X)\) is exact;
        \item The natural map \(\mc{O}_X\to q_*\mc{O}_{\mf{X}}\) is an isomorphism.
    \end{enumerate}
\end{defn}

\begin{table}[htpb]
    \centering
    \caption{Types of moduli spaces}
    \label{tab:moduli-types-corrected}
    \renewcommand{\arraystretch}{1.3}
    \begin{tabular}{
        p{0.18\textwidth}
        p{0.30\textwidth}
        p{0.22\textwidth}
        p{0.20\textwidth}
    }
        \toprule
        \textbf{Notion}
        & \textbf{Defining property}
        & \textbf{Geometric behavior}
        & \textbf{Feature}
        \\
        \midrule
        Fine moduli
        & Represents the set-valued moduli functor
        & Its points represent the objects recorded by the functor
        & Admits a universal family
        \\
        Coarse moduli
        & Universal for maps to algebraic spaces and bijective on geometric isomorphism classes
        & One geometric point for each isomorphism class
        & Forgets automorphisms
        \\
        Good moduli
        & Cohomologically affine and
          \(\mc{O}_X \simeq q_*\mc{O}_{\mf{X}}\)
        & May identify distinct nonclosed points
        & Gives \(S\)-equivalence in GIT-type examples
        \\
        \bottomrule
    \end{tabular}
\end{table}

The existence of these moduli spaces is not determined only by whether stabilizers are finite or infinite. For example, \(B\G_m\to\on{Spec}\C\) is both a coarse moduli space and a good moduli space, even though its stabilizer is \(\G_m\).

\begin{thm}[Keel-Mori {\cite{KeelMori1997}}]
    Let \(\mf{X}\) be an algebraic stack of finite type over \(\C\). If its inertia morphism
    \[
        I_{\mf{X}}
        \coloneqq
        \mf{X}\times_{\mf{X}\times\mf{X}}\mf{X}
        \longrightarrow \mf{X}
    \]
    is finite, then \(\mf{X}\) admits a coarse moduli space. The fiber of \(I_{\mf{X}}\to\mf{X}\) over a geometric point \(x\) is the automorphism group \(\on{Aut}(x)\).
\end{thm}

\begin{thm}[Alper--Hall--Rydh {\cite{AHR2020}}]
    Let
    \[
        q\colon \mf{X}\longrightarrow X
    \]
    be a good moduli space, where \(\mf{X}\) is a Noetherian
    algebraic stack over \(\C\) with affine diagonal.
    Let \(x\in\mf{X}(\C)\) be a closed point, and denote its
    stabilizer group by \(G_x\).

    Then there exists an affine scheme \(\on{Spec}A\) equipped
    with an action of \(G_x\), together with a Cartesian diagram
    \[
        \begin{tikzcd}
            {[\on{Spec}A/G_x]}
                \arrow[r]
                \arrow[d]
            & \mf{X}
                \arrow[d, "q"]
            \\
            \on{Spec}(A^{G_x})
                \arrow[r, "\mathrm{\acute{e}tale}"]
            & X,
        \end{tikzcd}
    \]
    where
    \[
        \on{Spec}(A^{G_x})\longrightarrow X
    \]
    is an \'etale neighborhood of \(q(x)\).
\end{thm}

For \(r\geq 2\), the stack \(\ms{Bun}_{r,d}(C)\) is locally of finite type, but it is neither quasi-compact nor separated. To see that it is not quasi-compact, choose line bundles \(L_n\) and \(M_{d-n}\) of degrees \(n\) and \(d-n\), respectively, and consider
\[
    E_n=L_n\oplus M_{d-n}\oplus\mc{O}_C^{\oplus r-2}.
\]
The bundle \(E_n\) contains a line subbundle of degree \(n\). In any quasi-compact family of vector bundles of fixed rank and degree, the degrees of line subbundles are bounded above, so the bundles \(E_n\) cannot all lie in a quasi-compact substack as \(n\to\infty\). To see that the stack is not separated, note that every vector bundle \(E\) has scalar automorphisms \(\G_m\subset\on{Aut}(E)\). Stabilizer groups of a separated algebraic stack are proper, whereas \(\G_m\) is not proper. Passing to semistable bundles does not make the stack separated; rather, it cuts out a bounded open substack of finite type which admits a good moduli space whose closed points are \(S\)-equivalence classes.

\begin{defn}[Slope stability]
    For a nonzero vector bundle \(E\), define its slope by
    \[ \mu(E)=\frac{\deg E}{\on{rk}E}. \]
    The vector bundle \(E\) is \textit{semistable} if
    \[ \mu(F)\leq\mu(E) \]
    for every nonzero proper saturated subsheaf \(F\subset E\), and it is \textit{stable} if the inequality is strict. Since \(C\) is a smooth curve, a saturated subsheaf of a vector bundle is a subbundle, so it is equivalent to test nonzero proper subbundles.
\end{defn}

Armed with this definition, we may now consider the open substack \(\ms{Bun}_{r,d}^{\mr{ss}}(C)\). Semistable bundles of fixed rank and degree form a bounded family, so this stack is of finite type, and in particular quasi-compact. It is still not separated because every object has the nonproper scalar stabilizer \(\G_m\).

\begin{defn}[S-equivalence]
    Let \(E\) be a semistable vector bundle. A Jordan--H\"older filtration of \(E\) is a filtration
    \[ 0=J_0\subset J_1\subset\cdots\subset J_N=E \]
    such that every quotient \(J_i/J_{i-1}\) is stable of slope \(\mu(E)\). The filtration is not necessarily unique, but the associated graded vector bundle
    \[
        \on{gr}(E)=\bigoplus_{i=1}^N J_i/J_{i-1}
    \]
    is unique up to noncanonical isomorphism. Two semistable vector bundles \(E\) and \(E'\) of the same slope are \textit{S-equivalent} if \(\on{gr}(E)\simeq\on{gr}(E')\).
\end{defn}

Now we consider the moduli space \(M_{r,d}^{\mr{ss}}(C)\) of \(S\)-equivalence classes of semistable vector bundles on \(C\). There is a good moduli space morphism
\[
    \ms{Bun}_{r,d}^{\mr{ss}}(C)\longrightarrow M_{r,d}^{\mr{ss}}(C),
\]
and the closed points of \(M_{r,d}^{\mr{ss}}(C)\) are in bijection with \(S\)-equivalence classes of semistable vector bundles.

\begin{rmk}
    The original method of proof for this theorem was to construct the moduli space using geometric invariant theory. The goal of this series of lectures will be to discuss an intrinsic approach using \(\Theta\)-stratifications together with existence criteria for good moduli spaces, without first choosing a GIT presentation.
\end{rmk}


